\documentclass[12pt,a4paper]{ctexart}

\usepackage{amsmath,amssymb,bm}
\usepackage{geometry}
\usepackage{enumitem}
\usepackage{booktabs}
\usepackage{xcolor}
\usepackage{hyperref}

\geometry{left=2.1cm,right=2.1cm,top=2.0cm,bottom=2.0cm}
\setlist[itemize]{itemsep=0.25em,topsep=0.25em}
\setlist[enumerate]{itemsep=0.35em,topsep=0.35em}

\newcommand{\dd}{\mathrm d}
\newcommand{\Rey}{\mathrm{Re}}
\newcommand{\blue}[1]{\textcolor{blue!70!black}{#1}}
\newcommand{\red}[1]{\textcolor{red!75!black}{#1}}

\title{第8章 粘性流动：作业补全笔记}
\author{}
\date{}

\begin{document}
\maketitle

\section*{0. 这次作业为什么你的原笔记不够}

你原来的笔记主要覆盖了“圆管层流”，也就是泊肃叶流动：
\[
u(r)=\frac{1}{4\mu}\left(-\frac{\dd p}{\dd x}\right)(R^2-r^2),
\qquad
\tau_w=\frac{R}{2}\left(-\frac{\dd p}{\dd x}\right),
\qquad
\lambda=\frac{64}{\Rey}.
\]

这能解决圆管压降、管内层流、阻力系数等题。
但这次作业还需要另外三块：

\begin{center}
\begin{tabular}{lll}
\toprule
知识块 & 典型题 & 核心公式 \\
\midrule
圆管层流 & 1、8.4、7 & \(\Rey=\dfrac{Vd}{\nu}\)，\(\tau_w=\dfrac{8\mu \bar V}{D}\) \\
小球低雷诺数 & 8.5、8.10 & \(F_D=6\pi\mu aU\) \\
平板剪切流 & 8.6、8 & \(\mu u''=\dfrac{\dd p}{\dd x}\)，\(\tau=\mu u'\) \\
涡量与流函数 & 9、10 & \(\Omega_z=v_x-u_y\)，\(u=\psi_y,\ v=-\psi_x\) \\
\bottomrule
\end{tabular}
\end{center}

\section*{1. 圆管层流：你原笔记的主线}

\subsection*{1.1 判断层流还是湍流}

圆管雷诺数：
\[
\boxed{
\Rey=\frac{Vd}{\nu}=\frac{\rho Vd}{\mu}
}
\]

常用判断：
\[
\Rey<2000\Rightarrow \text{层流},
\qquad
\Rey>4000\Rightarrow \text{湍流}.
\]

中间区域为过渡区，作业里一般按题目提示或常用界限判断。

\subsection*{1.2 圆管泊肃叶流动}

圆管半径为 \(R\)，直径为 \(D=2R\)。若流动为定常、充分发展、水平圆管层流，则
\[
\boxed{
u(r)=\frac{1}{4\mu}\left(-\frac{\dd p}{\dd x}\right)(R^2-r^2)
}
\]

最大速度：
\[
\boxed{
u_{\max}=\frac{R^2}{4\mu}\left(-\frac{\dd p}{\dd x}\right)
}
\]

平均速度：
\[
\boxed{
\bar V=\frac{u_{\max}}{2}
=\frac{R^2}{8\mu}\left(-\frac{\dd p}{\dd x}\right)
}
\]

流量：
\[
\boxed{
Q=\bar V\pi R^2
=\frac{\pi R^4}{8\mu}\left(-\frac{\dd p}{\dd x}\right)
}
\]

\subsection*{1.3 壁面剪应力与平均速度}

圆管层流中最常用的一条：
\[
\boxed{
\tau_w=\frac{8\mu \bar V}{D}
}
\]

所以如果已知 \(\tau_w,D,\bar V\)，求粘度：
\[
\boxed{
\mu=\frac{\tau_wD}{8\bar V}
}
\]

其中
\[
\bar V=\frac{Q}{A}=\frac{4Q}{\pi D^2}.
\]

\subsection*{1.4 圆管压降}

水平圆管：
\[
\boxed{
\tau_w=\frac{D}{4}\left(-\frac{\dd p}{\dd x}\right)
}
\]

所以
\[
\boxed{
-\frac{\dd p}{\dd x}=\frac{4\tau_w}{D}
}
\]

若是竖直管，还要考虑重力沿流动方向的分量。
取 \(x\) 沿流动方向，则
\[
\boxed{
\frac{\dd p}{\dd x}
=\rho g_x-\frac{4\tau_w}{D}
}
\]

其中：
\[
\begin{cases}
g_x=0, & \text{水平流动},\\
g_x=-g, & \text{向上流动},\\
g_x=+g, & \text{向下流动}.
\end{cases}
\]

\section*{2. 小球低雷诺数：补充你缺的第一块}

\subsection*{2.1 什么时候用}

看到“小球”“很慢”“粘度很大”“\(Re<1\)”时，优先想到斯托克斯阻力。

\[
\boxed{
\text{小球低雷诺数} \Rightarrow F_D=6\pi\mu aU
}
\]

其中 \(a\) 是小球半径，不是直径。

\subsection*{2.2 斯托克斯阻力}

小球半径 \(a\)，速度 \(U\)，流体粘度 \(\mu\)，阻力为
\[
\boxed{
F_D=6\pi\mu aU
}
\]

低雷诺数判断：
\[
\boxed{
\Rey=\frac{\rho Ud}{\mu}
=\frac{2\rho Ua}{\mu}<1
}
\]

\subsection*{2.3 小球匀速下落}

小球受三种力：
\[
\text{重力}=G=\frac43\pi a^3\rho_s g,
\]
\[
\text{浮力}=B=\frac43\pi a^3\rho g,
\]
\[
\text{阻力}=D=6\pi\mu aU.
\]

匀速时合力为零：
\[
G-B-D=0.
\]

因此
\[
\frac43\pi a^3(\rho_s-\rho)g=6\pi\mu aU.
\]

解得最终速度：
\[
\boxed{
U_t=\frac{2a^2(\rho_s-\rho)g}{9\mu}
}
\]

\subsection*{2.4 反过来求粘度}

若已知小球匀速下落速度 \(U\)，求流体粘度：
\[
\boxed{
\mu=
\frac{\frac43\pi a^3(\rho_s-\rho)g}{6\pi aU}
}
\]

如果题目直接给质量 \(m\)，不要先算 \(\rho_s\) 也可以：
\[
\boxed{
\mu=
\frac{mg-\rho gV_s}{6\pi aU},
\qquad
V_s=\frac43\pi a^3
}
\]

\section*{3. 平板剪切流：补充你缺的第二块}

\subsection*{3.1 基本方程}

平板间一维平行流：
\[
u=u(y),\qquad v=0.
\]

若定常、不可压、牛顿流体，则 \(x\) 方向动量方程常化为
\[
\boxed{
\mu\frac{\dd^2u}{\dd y^2}
=\frac{\dd p}{\dd x}-\rho g_x
}
\]

也常写成
\[
\boxed{
\mu u''-\frac{\dd p}{\dd x}+\rho g_x=0
}
\]

具体符号取决于 \(x\) 正方向怎么选，考试时只要前后一致。

\subsection*{3.2 无压力差、无体力：速度是直线}

如果
\[
\frac{\dd p}{\dd x}=0,\qquad g_x=0,
\]
则
\[
\mu u''=0.
\]

所以
\[
\boxed{
u=Ay+B
}
\]

这就是简单库埃特流。

\subsection*{3.3 两层流体夹在两平板之间}

下层厚度 \(h_1\)，粘度 \(\mu_1\)；
上层厚度 \(h_2\)，粘度 \(\mu_2\)。
下板静止，上板速度为 \(U\)。

关键不是从头解微分方程，而是记住：
\[
\boxed{
\text{两层剪应力相等}
}
\]

即
\[
\boxed{
\tau=\mu_1\frac{\dd u_1}{\dd y}
=\mu_2\frac{\dd u_2}{\dd y}
}
\]

每层速度差分别为
\[
\Delta u_1=\frac{\tau h_1}{\mu_1},
\qquad
\Delta u_2=\frac{\tau h_2}{\mu_2}.
\]

总速度差为
\[
U=\frac{\tau h_1}{\mu_1}+\frac{\tau h_2}{\mu_2}.
\]

所以
\[
\boxed{
\tau=
\frac{U}{h_1/\mu_1+h_2/\mu_2}
=
\frac{U\mu_1\mu_2}{\mu_2h_1+\mu_1h_2}
}
\]

取 \(y=0\) 为下板，\(y=h_1\) 为交界面，\(y=h_1+h_2\) 为上板，则
\[
\boxed{
u_1(y)=\frac{U\mu_2y}{\mu_2h_1+\mu_1h_2},
\qquad 0\le y\le h_1
}
\]

\[
\boxed{
u_2(y)=
\frac{U\left[\mu_2h_1+\mu_1(y-h_1)\right]}
{\mu_2h_1+\mu_1h_2},
\qquad h_1\le y\le h_1+h_2
}
\]

\subsection*{3.4 倾斜平板，有重力沿斜面分量}

取 \(x\) 沿斜面向下，\(y\) 从下板指向上板。

若沿流动方向无压力差：
\[
\frac{\dd p}{\dd x}=0,
\]
则
\[
\boxed{
\mu\frac{\dd^2u}{\dd y^2}+\rho g\sin\theta=0
}
\]

令
\[
A=\frac{\rho g\sin\theta}{\mu}.
\]

则
\[
u''=-A.
\]

积分：
\[
u=-\frac{A}{2}y^2+C_1y+C_2.
\]

如果下板不动：
\[
u(0)=0\Rightarrow C_2=0.
\]

如果中面速度为零：
\[
u(H/2)=0\Rightarrow C_1=\frac{AH}{4}.
\]

所以
\[
\boxed{
u=-\frac{A}{2}y^2+\frac{AH}{4}y
}
\]

若上板向上运动，与 \(x\) 正方向相反，则
\[
u(H)=-U_0.
\]

代入可得
\[
\boxed{
U_0=\frac{AH^2}{4}
=\frac{\rho gH^2\sin\theta}{4\mu}
}
\]

剪应力：
\[
\boxed{
\tau=\mu\frac{\dd u}{\dd y}
=\rho g\sin\theta\left(\frac H4-y\right)
}
\]

下壁面：
\[
\boxed{
|\tau_{\text{下}}|
=\frac{\rho gH\sin\theta}{4}
}
\]

上壁面：
\[
\boxed{
|\tau_{\text{上}}|
=\frac{3\rho gH\sin\theta}{4}
}
\]

\section*{4. 涡量与流函数证明题：补充你缺的第三块}

\subsection*{4.1 涡量定义}

三维涡量：
\[
\boxed{
\bm\Omega=\nabla\times\bm V
}
\]

二维 \(x-y\) 平面流动：
\[
\bm V=(u,v,0).
\]

此时涡量只有 \(z\) 方向：
\[
\boxed{
\Omega_z=\frac{\partial v}{\partial x}
-\frac{\partial u}{\partial y}
}
\]

\subsection*{4.2 静止壁面上的涡量法向分量}

静止固体壁面上满足无滑移条件：
\[
u=v=w=0.
\]

取壁面局部坐标：\(x,y\) 为切向，\(z\) 为法向。
则
\[
\bm n=\bm k.
\]

法向涡量为
\[
\bm\Omega\cdot\bm n=\Omega_z
=\frac{\partial v}{\partial x}
-\frac{\partial u}{\partial y}.
\]

因为壁面上 \(u,v\) 沿壁面处处为零，所以切向导数为零：
\[
\frac{\partial v}{\partial x}=0,
\qquad
\frac{\partial u}{\partial y}=0.
\]

故
\[
\boxed{
\bm\Omega\cdot\bm n=0
}
\]

\subsection*{4.3 流函数定义}

二维不可压缩流动可以引入流函数 \(\psi\)。
本笔记采用约定：
\[
\boxed{
u=\frac{\partial\psi}{\partial y},
\qquad
v=-\frac{\partial\psi}{\partial x}
}
\]

这个约定自动满足不可压连续方程：
\[
\frac{\partial u}{\partial x}
+\frac{\partial v}{\partial y}
=
\frac{\partial^2\psi}{\partial x\partial y}
-
\frac{\partial^2\psi}{\partial y\partial x}
=0.
\]

\subsection*{4.4 流函数与涡量的关系}

由
\[
\Omega_z=v_x-u_y,
\]
代入
\[
u=\psi_y,\qquad v=-\psi_x,
\]
得
\[
\Omega_z
=(-\psi_x)_x-(\psi_y)_y
=-\psi_{xx}-\psi_{yy}.
\]

所以
\[
\boxed{
\Omega=-\nabla^2\psi
}
\]

即
\[
\boxed{
\nabla^2\psi=-\Omega
}
\]

\section*{5. 作业题号对照表}

\begin{center}
\begin{tabular}{cll}
\toprule
题号 & 题型 & 你该用什么 \\
\midrule
1 & 判断层流/湍流 & \(\Rey=Vd/\nu\) \\
3：书8.4 & 圆管层流求粘度 & \(\tau_w=8\mu\bar V/D\) \\
4：书8.5 & 小球匀速下落求粘度 & \(mg-\rho gV_s=6\pi\mu aU\) \\
5：书8.6 & 两层平板剪切流 & \(\tau=\mu_1u_1'=\mu_2u_2'\) \\
6：书8.10 & 小球终速证明 & \(G-B-D=0\) \\
7 & 圆管层流加重力方向 & \(dp/dx=\rho g_x-4\tau_w/D\) \\
8 & 倾斜平板剪切流 & \(\mu u''+\rho g\sin\theta=0\) \\
9 & 静止壁面涡量证明 & 无滑移 \(u=v=0\Rightarrow\Omega\cdot n=0\) \\
10 & 流函数与涡量 & \(u=\psi_y,\ v=-\psi_x\Rightarrow\nabla^2\psi=-\Omega\) \\
\bottomrule
\end{tabular}
\end{center}

\section*{6. 最短背诵版}

\[
\boxed{
\Rey=\frac{Vd}{\nu}=\frac{\rho Vd}{\mu}
}
\]

\[
\boxed{
\tau_w=\frac{8\mu \bar V}{D},
\qquad
\tau_w=\frac{D}{4}\left(-\frac{\dd p}{\dd x}\right)
}
\]

\[
\boxed{
F_D=6\pi\mu aU
}
\]

\[
\boxed{
U_t=\frac{2a^2(\rho_s-\rho)g}{9\mu}
}
\]

\[
\boxed{
\tau=\mu\frac{\dd u}{\dd y}
}
\]

\[
\boxed{
\text{平板无压差无体力： }u''=0,\quad u=Ay+B
}
\]

\[
\boxed{
\text{倾斜平板： }\mu u''+\rho g\sin\theta=0
}
\]

\[
\boxed{
\Omega_z=v_x-u_y
}
\]

\[
\boxed{
u=\psi_y,\qquad v=-\psi_x,\qquad \nabla^2\psi=-\Omega
}
\]

\end{document}
